\worksheettitle{Наибольшее и наименьшее число}{\modulemark{M06} Учебный модуль \textbullet\ выбор по старшему разряду}
\studentnameblock

\begin{checkpointbox}[Вход после M05]
  Из цифр 0, 3, 7 составьте два трёхзначных числа без повторений:
  \answerblank[25mm], \answerblank[25mm]. Почему запись $037$ не подходит?
  \answerblank[75mm].
\end{checkpointbox}

\section{Как получить наибольшее число}
Из цифр 2, 5, 7, 9 нужно составить наибольшее четырёхзначное число, используя
каждую ровно один раз.

\begin{fillbox}[Выбираем слева направо]
  В разряд тысяч ставим наибольшую цифру: \answerblank[12mm].
  Затем выбираем наибольшую из оставшихся: \answerblank[12mm].
  Получаем число \answerblank[45mm].
\end{fillbox}

\begin{reflectionbox}[Почему способ работает]
  При сравнении чисел сначала смотрят на старший разряд. Любое число,
  начинающееся с 9, больше числа, начинающегося с 7, даже если дальше стоят
  большие цифры. Поэтому для максимума цифры располагают по убыванию.
\end{reflectionbox}

\begin{practicebox}[\levelbase. Наибольшее число]
  Используйте каждую цифру ровно один раз.
  \begin{enumerate}[label=\textbf{\arabic*.}]
    \item 1, 4, 6: \answerblank[35mm];
    \item 2, 3, 8, 9: \answerblank[35mm];
    \item 0, 4, 5, 7: \answerblank[35mm];
    \item 0, 0, 3, 8, 9: \answerblank[40mm].
  \end{enumerate}
\end{practicebox}

\section{Минимум без нуля}
Для цифр 2, 5, 7, 9 наименьшее число получается при расположении цифр по
\answerblank[30mm]. Число: \answerblank[38mm].

Объясните через разряд тысяч: \answerblank[65mm].

\newpage
\section{Минимум с нулём}
Из цифр 0, 2, 5, 7 нельзя получить четырёхзначное число $0257$. Поэтому
простого расположения по возрастанию недостаточно.

\begin{fillbox}[Точный алгоритм]
  \begin{enumerate}[label=\textbf{\arabic*.}]
    \item На первое место поставьте наименьшую \textbf{ненулевую} цифру.
    \item Сразу после неё поставьте все нули.
    \item Остальные цифры расположите по возрастанию.
  \end{enumerate}
  Для 0, 2, 5, 7 получаем \answerblank[45mm].
\end{fillbox}

\begin{warningbox}[Почему не $2507$]
  Числа $2057$ и $2507$ впервые различаются в разряде сотен:
  $0<5$, поэтому $2057<2507$. Ноль нужно поставить как можно левее, но не в
  начало записи.
\end{warningbox}

\begin{practicebox}[\levelbase. Максимум и минимум]
  \begin{center}
  \renewcommand{\arraystretch}{1.65}
  \begin{tabular}{c|c|c}
    \toprule
    цифры & наибольшее & наименьшее \\
    \midrule
    1, 3, 6, 8 & \answerblank[28mm] & \answerblank[28mm] \\
    0, 3, 6, 8 & \answerblank[28mm] & \answerblank[28mm] \\
    0, 1, 4, 9 & \answerblank[28mm] & \answerblank[28mm] \\
    0, 0, 5, 7 & \answerblank[28mm] & \answerblank[28mm] \\
    0, 2, 2, 6 & \answerblank[28mm] & \answerblank[28mm] \\
    \bottomrule
  \end{tabular}
  \end{center}
\end{practicebox}

\begin{warningbox}[Исправьте ошибку]
  Из цифр 0, 3, 5, 8 ученик составил минимум $3058$, но проверяющий объявил
  ответ неверным и предложил $3508$. Кто прав? \answerblank[35mm].
  Обоснуйте сравнением разряда сотен: \answerblank[90mm].
\end{warningbox}

\newpage
\section{Закрепление: выбирает учитель}

\begin{practicebox}[\levelpractice. Заполните таблицу]
  \begin{center}
  \renewcommand{\arraystretch}{1.65}
  \begin{tabular}{c|c|c}
    \toprule
    цифры & максимум & минимум \\
    \midrule
    3, 4, 7 & \answerblank[26mm] & \answerblank[26mm] \\
    1, 5, 5, 9 & \answerblank[26mm] & \answerblank[26mm] \\
    0, 2, 6, 9 & \answerblank[26mm] & \answerblank[26mm] \\
    0, 0, 4, 8, 9 & \answerblank[26mm] & \answerblank[26mm] \\
    0, 1, 1, 6, 7 & \answerblank[26mm] & \answerblank[26mm] \\
    2, 3, 4, 7, 8, 9 & \answerblank[26mm] & \answerblank[26mm] \\
    \bottomrule
  \end{tabular}
  \end{center}
\end{practicebox}

\begin{practicebox}[Объясните без полного перебора]
  \begin{enumerate}[label=\textbf{\arabic*.}]
    \item Почему из цифр 1, 4, 8 число $841$ больше всех остальных?
      \answerblank[72mm].
    \item Почему из цифр 0, 2, 7 число $207$ меньше $270$?
      \answerblank[72mm].
    \item Как изменится максимум из 2, 5, 9, если добавить цифру 0?
      \answerblank[70mm].
  \end{enumerate}
\end{practicebox}

\newpage
\section{Проверяем готовые решения}

\begin{warningbox}[Найдите разные ошибки]
  \begin{enumerate}[label=\textbf{\arabic*.}]
    \item Для цифр 2, 4, 7 максимум записали $7422$.
    \item Для цифр 0, 1, 6, 9 минимум записали $0169$.
    \item Для цифр 0, 3, 5, 8 минимум записали $3508$.
    \item Для цифр 1, 2, 7, 9 максимум записали $9271$.
  \end{enumerate}
  Исправьте ответы и подпишите причину каждой ошибки.
  \vspace{35mm}
\end{warningbox}

\begin{fillbox}[Восстановите набор]
  Наибольшее число из четырёх цифр равно $8531$. Какие цифры были даны?
  \answerblank[60mm]. Какое наименьшее число можно из них составить?
  \answerblank[40mm].

  Максимум равен $9700$. Набор цифр: \answerblank[55mm]. Минимум:
  \answerblank[40mm].
\end{fillbox}

\newpage
\section{Контрольная точка}
\begin{checkpointbox}[Выполните без подсказки]
  Из цифр 0, 2, 5, 8, используя каждую ровно один раз, составьте:
  \begin{enumerate}[label=\textbf{\arabic*.}]
    \item наибольшее число: \answerblank[40mm];
    \item наименьшее число: \answerblank[40mm];
    \item объясните расположение нуля и сравнение по старшему различающемуся
      разряду: \answerblank[115mm].
  \end{enumerate}
\end{checkpointbox}

\begin{reflectionbox}[Моя готовность]
  \choicebox\ Для максимума выбираю цифры по убыванию.

  \choicebox\ Для минимума не ставлю ноль первым.

  \choicebox\ Объясняю ответ через первый различающийся разряд.

  \choicebox\ Мне нужен пошаговый маршрут M06-R.
\end{reflectionbox}

\begin{notebox}[Если закончили раньше]
  Выполните M06\(\star\): найдите второе по величине число и обоснуйте, почему
  между ним и максимумом нет другого варианта.
\end{notebox}
